\documentclass[10pt,a4paper]{amsart}
\usepackage{amsmath,amssymb,amsthm,graphicx}
\usepackage{verbatim}
\newenvironment{reference}{\comment}{\endcomment}
%\newenvironment{reference}{}{}
\newenvironment{slogan}{\comment}{\endcomment}
\newenvironment{history}{\comment}{\endcomment}

% For commutative diagrams we use Xy-pic
\usepackage[all]{xy}

% We use 2cell for 2-commutative diagrams.
\xyoption{2cell}
% Theorem environments.
%
\theoremstyle{plain}
\newtheorem{theorem}[subsection]{Theorem}
\newtheorem{proposition}[subsection]{Proposition}
\newtheorem{lemma}[subsection]{Lemma}

\theoremstyle{definition}
\newtheorem{definition}[subsection]{Definition}
\newtheorem{example}[subsection]{Example}
\newtheorem{exercise}[subsection]{Exercise}
\newtheorem{situation}[subsection]{Situation}

\theoremstyle{remark}
\newtheorem{remark}[subsection]{Remark}
\newtheorem{remarks}[subsection]{Remarks}

\numberwithin{equation}{subsection}

% Macros
%
\def\lim{\mathop{\mathrm{lim}}\nolimits}
\def\colim{\mathop{\mathrm{colim}}\nolimits}
\def\Spec{\mathop{\mathrm{Spec}}}
\def\Hom{\mathop{\mathrm{Hom}}\nolimits}
\def\Ext{\mathop{\mathrm{Ext}}\nolimits}
\def\SheafHom{\mathop{\mathcal{H}\!\mathit{om}}\nolimits}
\def\SheafExt{\mathop{\mathcal{E}\!\mathit{xt}}\nolimits}
\def\Sch{\mathit{Sch}}
\def\Mor{\mathop{\mathrm{Mor}}\nolimits}
\def\Ob{\mathop{\mathrm{Ob}}\nolimits}
\def\Sh{\mathop{\mathit{Sh}}\nolimits}
\def\NL{\mathop{N\!L}\nolimits}
\def\CH{\mathop{\mathrm{CH}}\nolimits}
\def\proetale{{pro\text{-}\acute{e}tale}}
\def\etale{{\acute{e}tale}}
\def\QCoh{\mathit{QCoh}}
\def\Ker{\mathop{\mathrm{Ker}}}
\def\Im{\mathop{\mathrm{Im}}}
\def\Coker{\mathop{\mathrm{Coker}}}
\def\Coim{\mathop{\mathrm{Coim}}}

% Boxtimes
%
\DeclareMathSymbol{\boxtimes}{\mathbin}{AMSa}{"02}

%
% Macros for moduli stacks/spaces
%
\def\QCohstack{\mathcal{QC}\!\mathit{oh}}
\def\Cohstack{\mathcal{C}\!\mathit{oh}}
\def\Spacesstack{\mathcal{S}\!\mathit{paces}}
\def\Quotfunctor{\mathrm{Quot}}
\def\Hilbfunctor{\mathrm{Hilb}}
\def\Curvesstack{\mathcal{C}\!\mathit{urves}}
\def\Polarizedstack{\mathcal{P}\!\mathit{olarized}}
\def\Complexesstack{\mathcal{C}\!\mathit{omplexes}}
% \Pic is the operator that assigns to X its picard group, usage \Pic(X)
% \Picardstack_{X/B} denotes the Picard stack of X over B
% \Picardfunctor_{X/B} denotes the Picard functor of X over B
\def\Pic{\mathop{\mathrm{Pic}}\nolimits}
\def\Picardstack{\mathcal{P}\!\mathit{ic}}
\def\Picardfunctor{\mathrm{Pic}}
\def\Deformationcategory{\mathcal{D}\!\mathit{ef}}
\usepackage{cleveref}
\crefname{originallemma}{original lemma}{original lemmas}
\Crefname{originallemma}{Original lemma}{Original lemmas}
\makeatletter
\newlabel{duality-lemma-gorenstein-base-change}{{0C03}{source}{}{sourcepointer-1}{}}
\newlabel{duality-lemma-gorenstein-base-change@cref}{{[originallemma][0][]0C03}{[1][0][]source}{}{}{}}
\makeatother
\crefname{originallemma}{original lemma}{original lemmas}
\Crefname{originallemma}{Original lemma}{Original lemmas}
\makeatletter
\newlabel{duality-lemma-gorenstein-lci}{{0BVA}{source}{}{sourcepointer-2}{}}
\newlabel{duality-lemma-gorenstein-lci@cref}{{[originallemma][0][]0BVA}{[1][0][]source}{}{}{}}
\makeatother
\crefname{originallemma}{original lemma}{original lemmas}
\Crefname{originallemma}{Original lemma}{Original lemmas}
\makeatletter
\newlabel{dualizing-lemma-flat-under-gorenstein}{{0BJL}{source}{}{sourcepointer-3}{}}
\newlabel{dualizing-lemma-flat-under-gorenstein@cref}{{[originallemma][0][]0BJL}{[1][0][]source}{}{}{}}
\makeatother
\crefname{originalreference}{original reference}{original references}
\Crefname{originalreference}{Original reference}{Original references}
\makeatletter
\newlabel{lemma-nodal-lci}{{0C4E}{source}{}{sourcepointer-4}{}}
\newlabel{lemma-nodal-lci@cref}{{[originalreference][0][]0C4E}{[1][0][]source}{}{}{}}
\makeatother
\crefname{originallemma}{original lemma}{original lemmas}
\Crefname{originallemma}{Original lemma}{Original lemmas}
\makeatletter
\newlabel{varieties-lemma-geometrically-reduced}{{035X}{source}{}{sourcepointer-5}{}}
\newlabel{varieties-lemma-geometrically-reduced@cref}{{[originallemma][0][]035X}{[1][0][]source}{}{}{}}
\makeatother
\crefname{originallemma}{original lemma}{original lemmas}
\Crefname{originallemma}{Original lemma}{Original lemmas}
\makeatletter
\newlabel{varieties-lemma-perfect-reduced}{{020I}{source}{}{sourcepointer-6}{}}
\newlabel{varieties-lemma-perfect-reduced@cref}{{[originallemma][0][]020I}{[1][0][]source}{}{}{}}
\makeatother
\title{Nodal curves are Gorenstein and geometrically reduced}
\author{The Stacks Project Authors}
\date{stacks-project:curves.tex:local-snapshot-20261012; GNU FDL 1.2 or later}
\begin{document}
\maketitle
\noindent\textit{Editorial note.} All delivered mathematical text is the author's own, copied from the pinned Stacks Project chapter file curves.tex; no statement or proof step was rewritten, reordered or removed.\par\medskip
\noindent\textit{Editorial note.} The counted proof is the author's complete argument: the lci-Gorenstein chain, the alternative base-change-completion route, the geometric-reducedness reduction, and the reduced normal-form check are all retained verbatim.\par\medskip
\noindent\textit{Editorial note.} The label lemma-facts-about-nodal-curves names the counted statement, and the proof environment immediately following it in the source is delivered verbatim as its proof.\par\medskip
\noindent\textit{Editorial note.} The counted statement is carried by the author's own lemma environment in the chapter source.\par\medskip
\noindent\textit{Allowed standard prerequisites (editorial).}\begin{itemize}
\item Local complete intersection property of nodal curves. Same chapter, Lemma lemma-nodal-lci; Duality for Schemes, Lemmas duality-lemma-gorenstein-lci and duality-lemma-gorenstein-base-change. Nodal curves are lci, hence Gorenstein; alternatively Gorensteinness can be checked after base change to the algebraic closure.
\item Gorensteinness under completion and the local normal forms. Dualizing Complexes, Lemma dualizing-lemma-flat-under-gorenstein; Same chapter, discussion of the completed local rings k[[t]] and k[[x, y]]/(xy). Supports the alternative by-hand verification that the completed local rings of a nodal curve are Gorenstein.
\item Geometric reducedness. Varieties, Lemmas varieties-lemma-perfect-reduced and varieties-lemma-geometrically-reduced. Reduces geometric reducedness to reducedness of X\_\{overline\{k\}\}, where the completed normal forms are visibly reduced.
\end{itemize}\medskip
\ifdefined\hypertarget\hypertarget{sourcepointer-1}{}\fi
\noindent\textit{Source reference.} Original lemma 0C03: sources/stacks-chow/tags.txt, source lines 15073--15073. Prerequisite pointer invoked at proof line(s) 4774; it names a standard fact of a sibling chapter that the counted argument uses, and every step of the counted proof built on it is retained verbatim in the delivered body.\par\medskip
\ifdefined\hypertarget\hypertarget{sourcepointer-2}{}\fi
\noindent\textit{Source reference.} Original lemma 0BVA: sources/stacks-chow/tags.txt, source lines 14903--14903. Prerequisite pointer invoked at proof line(s) 4771; it names a standard fact of a sibling chapter that the counted argument uses, and every step of the counted proof built on it is retained verbatim in the delivered body.\par\medskip
\ifdefined\hypertarget\hypertarget{sourcepointer-3}{}\fi
\noindent\textit{Source reference.} Original lemma 0BJL: sources/stacks-chow/tags.txt, source lines 14529--14529. Prerequisite pointer invoked at proof line(s) 4777; it names a standard fact of a sibling chapter that the counted argument uses, and every step of the counted proof built on it is retained verbatim in the delivered body.\par\medskip
\ifdefined\hypertarget\hypertarget{sourcepointer-4}{}\fi
\noindent\textit{Source reference.} Original reference 0C4E: sources/stacks-curves/pinned-original-source.tex, source lines 4700--4700. Prerequisite pointer invoked at proof line(s) 4770; it names a standard fact of the same chapter that the counted argument uses, and every step of the counted proof built on it is retained verbatim in the delivered body.\par\medskip
\ifdefined\hypertarget\hypertarget{sourcepointer-5}{}\fi
\noindent\textit{Source reference.} Original lemma 035X: sources/stacks-chow/tags.txt, source lines 4144--4144. Prerequisite pointer invoked at proof line(s) 4785; it names a standard fact of a sibling chapter that the counted argument uses, and every step of the counted proof built on it is retained verbatim in the delivered body.\par\medskip
\ifdefined\hypertarget\hypertarget{sourcepointer-6}{}\fi
\noindent\textit{Source reference.} Original lemma 020I: sources/stacks-chow/tags.txt, source lines 2635--2635. Prerequisite pointer invoked at proof line(s) 4784; it names a standard fact of a sibling chapter that the counted argument uses, and every step of the counted proof built on it is retained verbatim in the delivered body.\par\medskip
\begin{lemma}
\label{lemma-facts-about-nodal-curves}
Let $k$ be a field. Let $X$ be a locally algebraic $k$-scheme
equidimensional of dimension $1$ whose singularities are at-worst-nodal.
Then $X$ is Gorenstein and geometrically reduced.
\end{lemma}

\begin{proof}
The Gorenstein assertion follows from Lemma \ref{lemma-nodal-lci}
and Duality for Schemes, Lemma \ref{duality-lemma-gorenstein-lci}.
Or you can use that it suffices to check after passing to the
algebraic closure (Duality for Schemes, Lemma
\ref{duality-lemma-gorenstein-base-change}), then use that
a Noetherian local ring is Gorenstein if and only if its
completion is so (by Dualizing Complexes, Lemma
\ref{dualizing-lemma-flat-under-gorenstein}), and
then prove that the local rings $k[[t]]$ and $k[[x, y]]/(xy)$
are Gorenstein by hand.

\medskip\noindent
To see that $X$ is geometrically reduced, it suffices to show that
$X_{\overline{k}}$ is reduced (Varieties, Lemmas
\ref{varieties-lemma-perfect-reduced} and
\ref{varieties-lemma-geometrically-reduced}).
But $X_{\overline{k}}$ is a nodal curve over an
algebraically closed field. Thus the complete local rings
of $X_{\overline{k}}$ are isomorphic to either
$\overline{k}[[t]]$ or $\overline{k}[[x, y]]/(xy)$
which are reduced as desired.
\end{proof}

\end{document}
